%% ============================================================
%% AULA 5 — Treinar é otimizar: descida de gradiente e variações
%% GBC073 — Inteligência Computacional (FACOM/UFU)
%%
%% Versão sucinta: roteiro do apresentador (poucas palavras,
%% mesmas definições e figuras do aula05.tex)
%%
%% Esqueleto com esboço de conteúdo — a desenvolver.
%% Ementa (ficha SEI 5116656): item 2 — aprendizagem supervisionada;
%% projeto de redes (parâmetros, modos de treinamento).
%%
%% Compilar:  xelatex aula05-sucinta.tex (duas passadas)
%% ============================================================
\documentclass[aspectratio=169, 11pt]{beamer}

\makeatletter
\def\input@path{{./}{../}}
\makeatother

\usetheme{InteligenciaComputacional}   % ANTES do babel: fontspec vem primeiro
\usepackage{babel}
\babelprovide[import, main]{brazilian}

\title[Aula 5 — Otimização]{Aula 5: Treinar é otimizar — descida de gradiente e suas variações}
\subtitle[GBC073 — Inteligência Computacional]{SGD e mini-batch \textbullet\ momentum \textbullet\ Adam \textbullet\ inicialização}
\author[Prof.\ Marcelo Albertini]{Prof.\ Marcelo Keese Albertini}
\institute{GBC073 — Inteligência Computacional \textbullet\ Universidade Federal de Uberlândia}
\date{2026}

\begin{document}

% ------------------------------------------------------------ capa
\begin{frame}[plain, noframenumbering]
  \titlepage
\end{frame}

% ------------------------------------------------------------ roteiro
\begin{frame}{Roteiro da aula}
  \begin{itemize}\setlength{\itemsep}{0.5ex}
    \item<1-> \textbf{As três descidas:} batch, estocástica e mini-batch
    \item<2-> \textbf{Momentum e descida acelerada}
    \item<3-> \textbf{Métodos adaptativos:} AdaGrad, RMSProp, Adam
    \item<4-> \textbf{Inicialização e taxas de aprendizado}
  \end{itemize}
\end{frame}

% ------------------------------------------------------------
\section{As três descidas}

\begin{frame}{As três descidas (batch, estocástica, mini-batch)}
  \begin{itemize}\setlength{\itemsep}{0.4ex}
    \item Em lote: gradiente exato, caro, sem ruído
    \item Estocástica (SGD): um exemplo por passo — ruído alto, escapa de
          mínimos ruins, converge devagar
    \item Mini-batch: o meio-termo prático; o ruído vira aliado
    \item Compromisso exploração\,/\,explotação dentro do próprio treino
  \end{itemize}
\end{frame}

% ------------------------------------------------------------
\section{Momentum e descida acelerada}

\begin{frame}{Momentum e descida acelerada}
  \begin{itemize}\setlength{\itemsep}{0.4ex}
    \item Momentum: acumula velocidade na direção consistente do gradiente
    \item Atravessa platôs e vales estreitos mais rápido
    \item NAG (Nesterov): avalia o gradiente \emph{à frente}, onde a
          velocidade já levou
    \item Interpretação física: bola rolando vs.\ degrau a cada passo
  \end{itemize}
\end{frame}

% ------------------------------------------------------------
\section{Métodos adaptativos}

\begin{frame}{Métodos adaptativos (AdaGrad, RMSProp, Adam)}
  \begin{itemize}\setlength{\itemsep}{0.4ex}
    \item AdaGrad: taxa por parâmetro, dividida pela soma histórica dos
          gradientes ao quadrado — morre devagar
    \item RMSProp: média móvel em vez de soma — esquece o passado distante
    \item Adam: momentum + RMSProp, com correção de viés no início do treino
    \item Adam é o padrão de fato em deep learning; entenda o que ajusta
  \end{itemize}
\end{frame}

% ------------------------------------------------------------
\section{Inicialização e taxas de aprendizado}

\begin{frame}{Inicialização e taxas de aprendizado}
  \begin{itemize}\setlength{\itemsep}{0.4ex}
    \item Inicialização importa: pesos constantes ou grandes demais matam o
          sinal (ou o gradiente)
    \item Xavier/Glorot e He: variância calibrada pela largura da camada
    \item Taxa de aprendizado: pequena demais é lenta; grande demais diverge
    \item Agendamentos (\emph{schedules}): decaimento, warm-up, cosine annealing
  \end{itemize}
  \begin{quizbox}
    Adam com taxa $\taxa = 1$ tende a convergir mais rápido do que com
    $\taxa = 10^{-3}$, pois aprende com passos maiores.

    \smallskip
    \opcoes{Verdadeiro \;/\; Falso}
  \end{quizbox}
\end{frame}

\end{document}
